
%% TODO: Centralize control over notations
% \renewcommand{\vec}[1]{\overrightarrow{\mathbf{#1}}}

%% TODO: Move notation here:
\section{Notations}

%% General mathematical notation
\begin{itemize}
    \item All logarithms are base 2, that is $\log x$ denotes $\log_2 x$.
    \item Vectors versus scalars etc.
    \item Vectors or sequences of elements are sometimes denoted using bracket notation. For example, a sequence of $k$ elements is written as $\{f(i)\}_{i=0}^{k-1}$.
    \item Both $\vec{v}[i]$ and $\vec{v}_i$ denotes the $i$-th element of $v$.
    \item For a vector $\vec{x}$ we use the max norm $\|\vec{x}\|_\infty$ to refer the largest element in the vector: $\max_i\{\vec{x}_i\}$ % Throughout, $\|\cdot\|$ denotes the infinity norm $\|\cdot\|_\infty$, that is, the largest absolute value among the coefficients of a polynomial (equivalently, among its residues in the representation of \autoref{sec:rns}).
\end{itemize}

%% Specialized notations
\begin{definition}[Centered Residue Representation] % TODO: Shorten
    \label{def:centered}
    Let $\mathbb{Z}_Q$ denote the quotient ring $\mathbb{Z}/Q\mathbb{Z}$. While $\mathbb{Z}_Q$ can be represented by any set of residue classes, in this work we canonically identify $\mathbb{Z}_Q$ with the set of centered residues in the interval $[-Q/2, Q/2)$. 
    
    We use the notation $[x]_Q$ to denote the unique representative of $x$ modulo $Q$ in this centered interval. When measuring the size of an element $a \in \mathbb{Z}_Q$ or a polynomial in $R_Q$, we inherently refer to the norm of its centered representative lifted to $\mathbb{Z}$ or $\mathbb{Z}[X]$.

    % From ealier
    For a modulus $q$, an element of $\mathbb{Z}_q$ is identified with its balanced representative through the map
    \begin{equation}
        [a]_q = \big( a + \lfloor q/2 \rfloor \bmod q \big) - \lfloor q/2 \rfloor .
    \end{equation}
    The flooring makes the map well defined for every $q$, and its image is the run of $q$ consecutive integers centered at the origin, so that $|[a]_q| \leq \lfloor q/2 \rfloor$. Working with balanced representatives is what allows every noise term to be bounded by $\|\epsilon\| \leq \lfloor q/2 \rfloor$ rather than by $q$.

    The notation extends to rings. If $S$ is a set of elements $\mathbb{Z}_Q^{N}$ and $a\in \mathbb{Z}_{Q}^{N}$ then $[a]_S = ([a_0]_{S_0}, \dots, [a_{N-1}]_{S_{N-1}})$, or if $a$ is a vector and $q$ is a scalar then $[a]_q$ applies to each coefficient.
\end{definition}

\begin{definition}[Working Ring]
    All plaintexts and ciphertexts are elements of the cyclotomic ring $R_Q = \mathbb{Z}_Q[X]/(X^N + 1)$, where $N$ is a power of two and $Q$ is the ciphertext modulus. Elements are polynomials of degree at most $N - 1$ whose coefficients are reduced into a centered range about zero.
    \begin{equation}
        \label{form:canonical}
        a = c_0 + c_1X + \cdots + c_{N-1}X^{N-1} = \sum_{i=0}^{N-1} c_iX^i
    \end{equation}
\end{definition}